\section{预备知识}\label{sec:prelim}

\subsection{Zeckendorf 语法与黄金均值子移位}
设二进制词 $w=w_1\cdots w_m\in\{0,1\}^m$。黄金均值语法约束为
\[
  w_iw_{i+1}=0 \quad (1\le i<m),
\]
即禁止子串 ``11''。令
\[
  X_m := \{w\in\{0,1\}^m:\ \text{$w$ 无相邻 $1$}\}.
\]
经典计数给出 $\abs{X_m}=F_{m+2}$（Fibonacci 数列）。

\subsection{图、覆盖与带标签覆盖}
我们使用有限有向图 $G=(V,E)$。覆盖映射 $\pi:\widetilde G\to G$ 的“局部同构”版本要求：对每个顶点 $v\in \widetilde G$，其出边/入边邻域与 $\pi(v)$ 的对应邻域等价；若在少数点上局部度结构改变，则为分支覆盖（branched cover）。若在边上附加标签，使得纤维信息可逆，则称为带标签覆盖。

\subsection{Ihara--Zeta 与 Bass 行列式公式}
对有限图 $G$ 的 Ihara zeta
\[
  Z_G(u)=\prod_{[p]}\bigl(1-u^{\ell(p)}\bigr)^{-1},
\]
其中 $[p]$ 遍历“素闭无回溯游走”等价类。Bass 行列式公式（亦可通过 Hashimoto 边算子表述）把 $Z_G(u)$ 写为邻接算子谱的行列式表达，从而建立 ``闭环计数 $\Leftrightarrow$ 谱'' 的确定性对应（见 \cite{bass1992ihara,ihara1966zeta}）。

\subsection{谱隙、Cheeger 与混合时间}
对 $d$-正则图的简单随机游走转移矩阵 $P=\tfrac1dA$，谱隙 $\gamma:=1-\lambda_2(P)$ 控制混合时间。标准结果给出混合时间对 $\gamma$ 的上/下界以及与 conductance 的 Cheeger 型夹逼（见 \cite{levinpereswilmer,lmspectralsgapnotes}）。

\subsection{Babai 图同构准多项式时间}
Babai 给出图同构（GI）准多项式时间算法（Johnson 图作为关键障碍结构），为复杂性端提供一个基准上界 \cite{babai2016gi}。

